\documentclass[12pt,a4paper]{article}
\usepackage{amsmath}
\usepackage{amssymb}  % Provides \mathbb and other symbols
\usepackage{graphicx}
\usepackage{subcaption}
\usepackage{enumitem}
\usepackage{listings}
\usepackage{hyperref}
\usepackage{ulem}
\usepackage{xcolor}
\usepackage{amsthm}
\usepackage{tcolorbox}     % For colored boxes
\usepackage[table]{xcolor} % For coloring table cells
\usepackage{array}         % For better table alignment
\usepackage{mdframed}
\usepackage{cancel}
\usepackage[margin=2.5cm]{geometry}
\usepackage{lmodern}
\setlength{\parindent}{0pt}
\setlength{\parskip}{0.8em}
\newcommand{\answer}{\vspace{3\baselineskip}\textbf{ANSWER:}}
\newcommand{\question}[2]{\textbf{Question #1}\quad #2}

\begin{document}

\begin{titlepage}
\centering
\vspace*{\fill} % Adds vertical space to center the title block vertically

{\Huge Complex Analysis} \\[0.5cm] % Title 2
{\Huge Home Work 2}\\[1.5cm] % Title 3
{\Large Student: Mark Roberts [26-712-224]}\\[0.5cm] % Student ID
{\Large \today} % Date, \today adds today's date

\vspace*{\fill}
\end{titlepage}
\newpage

\question{2.1}{Look again at the subsets in Exercise 1.2(a) from Problem Sheet 1. Which of these sets are domains?}

\vspace{1cm}

\textbf{ANSWER:}\\
\\
A domain (aka region) is a connected open subset of $\mathbb{C}$.\\
\\
The case (a.i) is NOT a domain.  It is connected but not open.\\
\\
The case (a.ii) is NOT a domain.  It is connected but not open.\\
\\
The case (a.iii) is NOT a domain.  It is open but not connected.\\
\\
The case (a.iv) is NOT a domain.  It is connected but not open.
\newpage

\question{2.2}{At which points $z \in \mathbb{C}$ are the following functions
(a) continuous? (b) complex-differentiable?}

\begin{enumerate}[label=(\roman*)]
    \item $\displaystyle \frac{1}{z^2+9}$

    \item $\displaystyle z^{10}+4\pi z^3-11$

    \item
    \[
    \begin{cases}
    \displaystyle \frac{z^2-1}{z^3+1}, & z\neq -1,\\[4pt]
    \frac{-2}{3}, & z=-1
    \end{cases}
    \]

    \item $(\operatorname{Re} z)(\operatorname{Im} z)$
\end{enumerate}

\vspace{1cm}

\textbf{ANSWER:}\\
(i)(a) Here we have a rational function - which is continuous at all points where the denominator is not zero.\\
The denominator is zero when: $z^2+9=0$, i.e. where $z = \pm 3i$.  Hence, $\frac{1}{z^2 + 9}$ is continuous 
at all points of $z \in \mathbb{C} \backslash \{-3i, 3i \}$.\\
\\
(b) Being a rational function it is also differentiable for $z \in \mathbb{C} \backslash \{-3i, 3i \}$.\\
\\
\\
(ii)(a) Here we have a standard polynomial function and hence this function is continuous for all $z \in \mathbb{C}$.\\
\\
(b) This function is also complex-differentiable at all $z \in \mathbb{C}$.\\
\\
\\
(iii)(a) For this function it is rational and continuous, except possibly at $z=-1$.  It is definitely \textbf{discontinuous} as:
$z=\frac{1 \pm i \sqrt{3}}{2}$ - the other places denominator is zero.
To check the limit, at $z \to -1$, let $z=-1+\epsilon$.  Then:
$$ z^2-1 = (-1 + \epsilon)^2 -1 = -2\epsilon + \epsilon^2$$.
Also:
$$z^3 + 1 = (-1 + \epsilon)^3 +1 = 3\epsilon - 3\epsilon^2 + \epsilon^3$$
Hence:
\begin{align*}
    \lim_{z \to -1}\frac{z^2-1}{z^3+1} &= \lim_{\epsilon \to 0}\frac{-2\epsilon + \epsilon^2}{3\epsilon - 3\epsilon^2 + \epsilon^3} \\
&= \lim_{\epsilon \to 0}\frac{-2 + \epsilon}{3 - 3\epsilon + \epsilon^2} \\
&= -\frac{2}{3}
\end{align*}
But this is the value of the function at $z=-1$ and hence the function is continuous on $\mathbb{C}$.

\newpage
\noindent
(b) Now: $f(z) = \frac{z^2-1}{z^3+1}$ has derivative:
\begin{align*}
    f'(z) &= (z^2-1)' (z^3+1)^{-1} + (z^2-1)[(z^3+1)^{-1}]' \\
    &= 2z(z^3+1)^{-1} + (z^2-1)[3z^2 (-1)](z^3+1)^{-2} \\
    &= \frac{1}{(z^3+1)^2} \left[2z(z^3+1) - 3z^2(z^2-1)) \right] \\
    &= \frac{-z^4+3z^2+2z}{(z^3+1)^2} \\
    &= \frac{z(2-z)(z+1)^2}{(z+1)^2(z^2-z+1)^2} \\
    &= \frac{z(2-z)}{(z^2-z+1)^2} \tag{2.2.1}
\end{align*}
and this will definitiely be valid for $z \in \mathbb{C} \setminus \{1, \frac{1 - i \sqrt{3}}{2}, \frac{1 + i \sqrt{3}}{2}\}$.\\
\\
As $f(z)$ is not continuous at $z=\frac{1 \pm i \sqrt{3}}{2}$ it can't be differentiable at these two points.\\
At $z=-1$ we see from (2.2.1) that $f'(z)$ exists and equals to $-1/3$.\\
\\
(iv)(a) Here the function can be represented by: $f(z)=xy$, where $z=x+iy$.\\
This is a polynomial in x and y and is therefore continuous on $\mathbb{C}$.\\
\\
(b) Here we will use the CR equations, where: $f(z) = u(x,y)+ iv(x,y)$ with: $u(x,y)=xy$ and $v(x,y)=0$.\\
\\
Clearly:
\begin{align*}
    u_x(x,y) &= y \quad u_y(x,y) = x \\
    v_x(x,y) &= 0 \quad v_y(x,y) = 0
\end{align*}
The differentials are continuous everywhere but $u_x = v_y$ only at $y=0$.  Also, $u_y = -v_x$ only at $x=0$.
consequently, the Cauchy Riemann equations are only satisfied at $(0,0)$ and hence we know that this function can, at most, be
differentiable at the origin.  As the partial derivatives are indeed continuous on the whole of the complex plane we can
conclude that the derivative does exist at the origin and the derivative is 0.


\newpage

\question{2.3}{
\begin{enumerate}[label=(\alph*)]
    \item Show that the function
    \[
    z \mapsto |z|^2
    \]
    is complex-differentiable only at the point $z=0$.

    \item Let $U \subset \mathbb{C}$ be an open set. Suppose that
    $f:U\to\mathbb{C}$ is complex-differentiable at some point
    $z\in U$. Must $f$ be holomorphic at $z$?

    \item At which points $z\in\mathbb{C}$ are the following functions
    (a) complex-differentiable, (b) holomorphic?
    
    \begin{enumerate}[label=(\roman*)]
        \item $|z|$
        \item $z|z|$
    \end{enumerate}
\end{enumerate}
}

\vspace{1cm}

\textbf{ANSWER:} \\
(a) Given $f(z) = |z|^2 = x^2 + y^2 = u(x,y) + i v(x,y)$ - the associated Cauchy Riemann (CR) equations are:
\begin{align*}
    u_x(x,y) &= 2x \quad u_y = 2y \\
    v_x(x,y) &= 0 \quad v_y = 0 \\
\end{align*}
Now, it is clear that $u_x = v_y$ only at $x=0$.  Also, it is clear that $u_y = -v_x$ only at $y=0$.\\
Hence, only at the origin (i.e. $z=0$) are the CR equations satisfied.  Additionally, the derivatives are continuous at the origin
and hence $f$ is differentiable at the origin with derivative 0.\\
\\
(b) To say that $f$ is holomorphic at a point $z_0 \in \mathbb{C}$ means that it is holomorphic at some open neighbourhood of $z_0$.
But if $f$ is differentiable at a point it does NOT mean that it is holomorphic at $z_0$.
Take the case (a) above, where $f$ is defined on (the open set) $\mathbb{C}$ but is only differentiable at one point - the origin.\\
\\
(c) (i) Here we have: $f(z) = \sqrt(x^2 + y^2) = u(x,y) + i v(x,y)$.  Then the CR equations give:
\begin{align*}
    u_x(x,y) &= \frac{x}{(x^2 + y^2)^{1/2}} \quad u_y(x,y) = \frac{y}{(x^2 + y^2)^{1/2}} \\
    v_x(x,y) &= 0 \quad\quad\quad\quad\quad\quad v_y(x,y) = 0
\end{align*}
The CR equations give us that $f$ is differentiable nowhere except possibly at the origin (i.e. at $z=0$).  However, $u_x$ and $u_y$ are not continuous
at the origin and hence can't be differentiable there either.\\
\\
It follows that $f$ is holomorphic nowhere.\\

\newpage
\noindent
(ii) Here we have: $f(z) = z|z| = (x + iy) \sqrt(x^2 + y^2) = u(x,y) + iv(x,y)$.\\
Hence: $u(x,y)=x \sqrt(x^2 + y^2)$ and $v(x,y)=y \sqrt(x^2 + y^2)$. From which we obtain the CR equations:
\begin{align*}
    u_x(x,y) &= \frac{2x^2 + y^2}{\sqrt{x^2 + y^2}} \quad u_y(x,y) = \frac{xy}{\sqrt{x^2 + y^2}} \\
    v_x(x,y) &= \frac{xy}{\sqrt{x^2 + y^2}} \quad v_y(x,y) = \frac{x^2 + 2y^2}{\sqrt{x^2 + y^2}}
\end{align*}
The CR equations: $u_x = v_y => x^2=y^2$ and $u_y = -u_x => xy=0$.  The only possibility for a solution would be at the origin.
Now all 4 equations are continuous at the origin - taking the value 0.  Hence $f$ is differentiable only at the origin and takes the
value 0.

\newpage

\question{2.4}{
Let $f:\mathbb{C}\to\mathbb{C}$ be the function
\[
f(z)=z^3.
\]

Show that there is no point $w$ on the line segment connecting $1$ and $i$, i.e.
\[
w\in\{ti+(1-t)1\mid t\in\mathbb{R},\ 0\leq t\leq 1\},
\]
such that
\[
\frac{f(i)-f(1)}{i-1}=f'(w).
\]
}

\vspace{1cm}

\textbf{ANSWER:} \\
\\
Taking $f(z) = z^3$ we determine:
\begin{align*}
\frac{f(i)-f(1)}{i-1} &= \frac{-i-1}{i-1} \\
 &= \frac{1+i}{1-i} \\
 &= \frac{1+i}{1-i} \frac{1+i}{1+i} \\
 &= i \tag{2.4.1}
\end{align*}

Now, as $f'(z) = 3z^2$, we need to find a $w \in \{(1-t) + ti \; | \; 0 \le t \le 1 \}$ such that:
\begin{align*}
\frac{f(i)-f(1)}{i-1} &= 3w^2 \tag{2.4.2}
\end{align*}
Using (2.4.1) and (2.4.2) and that they must be equal, we obtain:
\begin{align*}
 3w^2 &= i \\
 \implies 3[(1-t) + ti]^2 &= i \\
 \implies 3\{(1-t)^2+2t(t-1)i-t^2\} &=i \\
 \implies (1-2t+t^2)+2t(t-1)i-t^2 &= i \\
 \implies (1-2t)+2t(t-1)i &= i.
\end{align*}
Real part:
\begin{align*}
  (1-2t) &= 0 \\
  \implies t &= 1/2
\end{align*}
Imaginary part:
\begin{align*}
  2t(t-1) &= 1 \qquad \text{no }t\in[0,1]\text{ solves this equation.}
\end{align*}

\begin{align*}
\therefore\ \text{no }w\in\{(1-t)+ti\mid 0\leq t\leq1\}\text{ can be found }
\end{align*}


\newpage

\question{2.5}{
Let $\Omega$ be a domain, and suppose that $z \in \Omega$. Show that
\[
\Omega \setminus \{z\}
\]
is a domain.

Now suppose that $r>0$ is such that
\[
D_r(z)\subseteq \Omega.
\]

Need
\[
\Omega \setminus D_r(z)
\]
be a domain?
}

\vspace{1cm}

\textbf{ANSWER:}\\
\\
If $\Omega \subset \mathbb{C}$ is a region/domain, then it is both open and connected.  Clearly, if $z \in \Omega$, then $\Omega \setminus \{z\}$ will still be open.\\
\\

\newpage

\question{2.6}{
Given $c \in \mathbb{C}$, let the sequence $z_n(c)$ be defined iteratively in the following way:
\[
z_0(c)=0,
\qquad
z_n(c)=z_{n-1}(c)^2+c
\quad \text{for } n\geq 1.
\]

The \textit{Mandelbrot set} $M$ is the set of $c\in\mathbb{C}$ such that $|c|\leq 2$ and $|z_n(c)|$ does not tend to infinity as $n\to\infty$.

\begin{enumerate}[label=(\roman*)]
    \item Suppose that there exists $k\in\mathbb{N}$ such that
    \[
    |z_k(c)|>2.
    \]
    Show that $c\notin M$.

    \item Let $U_k$ denote the set
    \[
    U_k:=\{c\in\mathbb{C}\mid |z_k(c)|>2\}.
    \]
    Show that $U_k$ is open.

    \item Deduce that $M$ is a closed subset of $\mathbb{C}$.
\end{enumerate}
}

\vspace{1cm}

\textbf{ANSWER:}

\end{document}
